% Figure 1, panel A: the data. Deliberately minimal; the experimental design of
% the m16h dataset belongs in the real-data section, not here. The column groups
% introduce the three node types, and their order (TFs, genes, enhancers) is the
% order used by the matrices in panels B and E. Everything else is caption work.
% Local origin is the panel's top-left corner. Panel size 5.3 x 4.8.
\draw[panel] (0,0) rectangle (5.3,-4.8);
\node[ttl] at (0.15,-0.12) {A\quad Single-cell multiome data};

% Shift measured with the frame and the panel letter excluded, so the
% drawing sits centred in the space below the letter rather than wherever
% its coordinates happened to put it.
\begin{scope}[shift={(+0.148,+0.630)}]  % centre the panel contents


\begin{scope}[shift={(1.20,-1.75)}]
  \def\cw{0.32}
  \def\ch{0.22}
  \pgfmathsetseed{7}
  \foreach \r in {1,...,14}{
    \foreach \c in {1,...,9}{
      \pgfmathsetmacro{\v}{4+78*rnd}
      \fill[black!\v!white, draw=white, line width=0.25pt]
        ({(\c-1)*\cw},{-(\r-1)*\ch}) rectangle ++(\cw,-\ch);
    }
  }

  % node-type bands above the columns, in the order used throughout the figure
  \foreach \a/\b/\col/\lab in {1/3/nodeTF/TFs, 4/6/nodeGene/genes, 7/9/nodeEnh/enhancers}{
    \draw[fill=\col!75, draw=none]
      ({(\a-1)*\cw+0.02},0.09) rectangle ({\b*\cw-0.02},0.17);
    \node[note, anchor=south] at ({(\a+\b-1)*\cw/2},0.19) {\lab};
  }

  \node[note, rotate=90, anchor=south] at (-0.16,-1.54) {$N$ cells};
  \node[note, anchor=north] at (1.44,-3.20) {$m$ nodes};
\end{scope}
\end{scope}
